% stephane [DOT] adjemian [AT] univ [DASH] lemans [DOT] fr
\documentclass[11pt,a4paper,notitlepage]{article}
\usepackage{amsmath}
\usepackage{amssymb}
\usepackage{amsbsy}
\usepackage{bm}
\usepackage[T1]{fontenc}
\usepackage[utf8x]{inputenc}
\usepackage{palatino}
\usepackage{scrtime}
\usepackage[french]{babel}
\usepackage{float}
\usepackage{nicefrac}

%%%%%%%%%%%%%%%%%%%%%%%%%%%%%%%%%%%%%%%%%%%%%%%%%%%%%%%%%%%%%%%%%%%%%%%%%%%%%%%%%%%%%%%%%%%%%%%%%%%%%

\newcommand{\exercice}[1]{\textsc{\textbf{Exercice}} #1}
\newcommand{\question}[1]{\textbf{(#1)}}
\setlength{\parindent}{0cm}



\begin{document}


\title{\textsc{Séries temporelles}\\(Éléments de correction)}
\author{\textsc{Université du Mans (Examen, L3)}}
\date{}


\maketitle

\exercice{1}\newline

\question{1} Le polynôme retard auto-régressif est~:
\[
\Phi(L) = 1 - \frac{3}{4}L + \frac{1}{8}L^2.
\]
Cherchons ses racines~: $\frac{1}{8}L^2 - \frac{3}{4}L + 1 = 0$, soit
$L^2 - 6L + 8 = 0$, dont les racines sont $L_1 = 2$ et $L_2 = 4$.
Toutes deux sont strictement supérieures à 1 en module, le processus
est donc asymptotiquement stationnaire au second ordre. On a
d'ailleurs la factorisation~:
\[
\Phi(L) = \left(1-\frac{1}{2}L\right)\left(1-\frac{1}{4}L\right).
\]
\newline

\question{2} Sous stationnarité, $\mathbb E[Y_t] = \mu$ pour tout $t$,
et l'application de l'espérance à l'équation du processus donne~:
\[
\mu = \frac{3/8}{1-3/4+1/8} = \frac{3/8}{3/8} = 1.
\]
\newline

\question{3} Notons $\tilde Y_t = Y_t - \mu$, qui vérifie~:
\[
\tilde Y_t = \frac{3}{4}\tilde Y_{t-1} - \frac{1}{8}\tilde Y_{t-2} + \varepsilon_t.
\]
En multipliant cette équation par $\tilde Y_{t-h}$ et en prenant
l'espérance, on obtient pour tout $h\geq 1$ les équations de
Yule-Walker~:
\[
\gamma(h) = \frac{3}{4}\gamma(h-1) - \frac{1}{8}\gamma(h-2),\quad h\geq 1
\]
(en utilisant $\gamma(-h)=\gamma(h)$). Pour $h=0$, on obtient en
outre, en multipliant par $\tilde Y_t$~:
\[
\gamma(0) = \frac{3}{4}\gamma(1) - \frac{1}{8}\gamma(2) + \sigma_{\varepsilon}^2.
\]

\textbf{Pour $h=1$~:} $\gamma(1) = \frac{3}{4}\gamma(0) -
\frac{1}{8}\gamma(1)$, donc $\frac{9}{8}\gamma(1) =
\frac{3}{4}\gamma(0)$, d'où~:
\[
\gamma(1) = \frac{2}{3}\gamma(0).
\]

\textbf{Pour $h=2$~:} $\gamma(2) = \frac{3}{4}\gamma(1) -
\frac{1}{8}\gamma(0) = \frac{3}{4}\cdot\frac{2}{3}\gamma(0) -
\frac{1}{8}\gamma(0) = \frac{3}{8}\gamma(0)$, soit~:
\[
\gamma(2) = \frac{3}{8}\gamma(0).
\]

\textbf{Pour $h=0$~:} en substituant ces deux relations,
\[
\gamma(0) = \frac{3}{4}\cdot\frac{2}{3}\gamma(0) - \frac{1}{8}\cdot\frac{3}{8}\gamma(0) + \sigma_{\varepsilon}^2,
\]
soit $\gamma(0)\left(1 - \frac{1}{2} + \frac{3}{64}\right) =
\sigma_{\varepsilon}^2$, c'est-à-dire $\frac{35}{64}\gamma(0) =
\sigma_{\varepsilon}^2$. Finalement~:
\[
\gamma(0) = \frac{64}{35}\sigma_{\varepsilon}^2,\quad
\gamma(1) = \frac{128}{105}\sigma_{\varepsilon}^2,\quad
\gamma(2) = \frac{24}{35}\sigma_{\varepsilon}^2.
\]
\newline

\question{4} La relation de récurrence est, pour tout $h\geq 1$~:
\[
\gamma(h) = \frac{3}{4}\gamma(h-1) - \frac{1}{8}\gamma(h-2).
\]
L'équation caractéristique associée est $r^2 - \frac{3}{4}r +
\frac{1}{8} = 0$, dont les racines sont $r_1=\frac{1}{2}$ et
$r_2=\frac{1}{4}$ (inverses des racines $L_1, L_2$ trouvées à la
question 1). La solution générale est donc~:
\[
\gamma(h) = A\left(\frac{1}{2}\right)^h + B\left(\frac{1}{4}\right)^h,
\]
les constantes $A$ et $B$ étant déterminées par les conditions
initiales $\gamma(0) = \frac{64}{35}\sigma_{\varepsilon}^2$ et
$\gamma(1) = \frac{128}{105}\sigma_{\varepsilon}^2$. La résolution du
système linéaire $A+B=\frac{64}{35}$, $\frac{A}{2}+\frac{B}{4} =
\frac{128}{105}$ donne $A = \frac{64}{21}\sigma_{\varepsilon}^2$ et $B
= -\frac{128}{105}\sigma_{\varepsilon}^2$. On a donc finalement~:
\[
\gamma(h) = \frac{64}{21}\left(\frac{1}{2}\right)^h\sigma_{\varepsilon}^2
- \frac{128}{105}\left(\frac{1}{4}\right)^h\sigma_{\varepsilon}^2.
\]
On note que $\gamma(h)$ tend vers 0 lorsque $h\to\infty$, ce qui est
cohérent avec la stationnarité du processus.

\bigskip

\exercice{2}\newline

\question{1} Le polynôme retard auto-régressif est $1+\frac{1}{2}L$,
dont la racine est $L=-2$, supérieure à 1 en module~: le processus
est asymptotiquement stationnaire au second ordre. Le polynôme retard
moyenne mobile est $1+\frac{1}{4}L$, dont la racine est $L=-4$,
également supérieure à 1 en module~: le processus est inversible. Ces
deux racines étant distinctes, on travaille bien sur la représentation
minimale d'un ARMA(1,1).\newline

\question{2} Sous stationnarité au second ordre, l'espérance et la
fonction d'autocovariance sont invariantes dans le temps~: $\mathbb
E[y_t]=\mu$ et $\mathbb E[(y_t-\mu)(y_{t-h}-\mu)]=\gamma(h)$ ne
dépendent pas de $t$.\newline

\question{3} En appliquant l'espérance à l'équation du processus~:
\[
\mu = 1 - \frac{1}{2}\mu \Longrightarrow \mu = \frac{2}{3}.
\]
\newline

\question{4} Posons $\tilde y_t = y_t - \mu$. Alors~:
\[
\tilde y_t = -\frac{1}{2}\tilde y_{t-1} + \varepsilon_t + \frac{1}{4}\varepsilon_{t-1}.
\]
On a besoin de $\mathbb E[\tilde y_t \varepsilon_t]$ et $\mathbb
E[\tilde y_t \varepsilon_{t-1}]$. Comme $\tilde y_t$ est, par
inversion du polynôme retard auto-régressif, une combinaison linéaire
des innovations courante et passées, on obtient directement~:
\[
\mathbb E[\tilde y_t \varepsilon_t] = \sigma_{\varepsilon}^2 = 1.
\]
Pour $\mathbb E[\tilde y_t \varepsilon_{t-1}]$, on substitue~:
\[
\begin{split}
\mathbb E[\tilde y_t \varepsilon_{t-1}] &= -\frac{1}{2}\mathbb E[\tilde y_{t-1}\varepsilon_{t-1}] + 0 + \frac{1}{4}\sigma_{\varepsilon}^2\\
&= -\frac{1}{2} + \frac{1}{4} = -\frac{1}{4}.
\end{split}
\]

En multipliant l'équation de $\tilde y_t$ par $\tilde y_t$ puis par
$\tilde y_{t-1}$ et en prenant les espérances~:
\[
\gamma(0) = -\frac{1}{2}\gamma(1) + 1 + \frac{1}{4}\cdot\left(-\frac{1}{4}\right)
= -\frac{1}{2}\gamma(1) + \frac{15}{16},
\]
\[
\gamma(1) = -\frac{1}{2}\gamma(0) + 0 + \frac{1}{4}\cdot 1
= -\frac{1}{2}\gamma(0) + \frac{1}{4}.
\]
En substituant la seconde dans la première, on obtient~:
\[
\gamma(0) = -\frac{1}{2}\left(-\frac{1}{2}\gamma(0)+\frac{1}{4}\right) + \frac{15}{16}
= \frac{1}{4}\gamma(0) + \frac{13}{16},
\]
soit $\frac{3}{4}\gamma(0) = \frac{13}{16}$, d'où~:
\[
\gamma(0) = \frac{13}{12},\quad
\gamma(1) = -\frac{1}{2}\cdot\frac{13}{12} + \frac{1}{4} = -\frac{7}{24}.
\]
\newline

\question{5} Pour $h\geq 2$, l'innovation $\varepsilon_{t-1}$ n'est
plus corrélée avec $\tilde y_{t-h}$ (qui ne dépend que de
$\varepsilon_{t-h},\varepsilon_{t-h-1},\dots$). En multipliant
l'équation de $\tilde y_t$ par $\tilde y_{t-h}$ et en prenant
l'espérance, il vient~:
\[
\gamma(h) = -\frac{1}{2}\gamma(h-1),\quad h\geq 2.
\]
Par itération, en partant de $\gamma(1)=-\frac{7}{24}$~:
\[
\gamma(h) = \left(-\frac{1}{2}\right)^{h-1}\gamma(1)
= -\frac{7}{24}\left(-\frac{1}{2}\right)^{h-1}, \quad h\geq 1.
\]
En particulier $\gamma(2) = \frac{7}{48}$, $\gamma(3) = -\frac{7}{96}$, etc.\newline

\question{6} L'autocorrélation est $\rho(h) =
\gamma(h)/\gamma(0)$. Comme $\gamma(0)>0$, le signe de $\rho(h)$ est
celui de $\gamma(h)$. La récurrence $\gamma(h) = -\frac{1}{2}\gamma(h-1)$
implique que $\gamma(h)$ change de signe à chaque pas pour $h\geq 1$.
Comme par ailleurs $\gamma(1)<0$, on a $\rho(h)>0$ pour $h$ pair et
$\rho(h)<0$ pour $h$ impair (sans oublier $\rho(0)=1>0$). La fonction
d'autocorrélation est donc \textit{alternée}, avec une décroissance
géométrique en module de raison $\frac{1}{2}$.

\bigskip

\exercice{3}\newline

\question{1} Notons $f$ la densité jointe de l'échantillon. En
appliquant le théorème de Bayes de manière itérée~:
\[
\begin{split}
f(\mathcal Y_T;\bm\theta) &= f(y_1,y_2;\bm\theta)\\
&\quad\times \prod_{t=3}^T f(y_t \mid y_{t-1},y_{t-2},\dots,y_1;\bm\theta).
\end{split}
\]
Or, dans le cas d'un AR(2), $y_t$ ne dépend que de
$(y_{t-1},y_{t-2})$ une fois conditionné par le passé, et l'innovation
$\varepsilon_t$ est gaussienne~: la loi conditionnelle est
$y_t\mid y_{t-1},y_{t-2} \sim \mathcal N(c+\varphi_1
y_{t-1}+\varphi_2 y_{t-2},\,\sigma_{\varepsilon}^2)$. La vraisemblance
exacte s'écrit donc~:
\[
\begin{split}
\mathcal L(\bm\theta;\mathcal Y_T) = & f(y_1,y_2;\bm\theta)\\
& \times \prod_{t=3}^T \frac{1}{\sqrt{2\pi\sigma_{\varepsilon}^2}}
e^{-\frac{(y_t - c - \varphi_1 y_{t-1} - \varphi_2 y_{t-2})^2}{2\sigma_{\varepsilon}^2}}.
\end{split}
\]

\textit{Difficulté par rapport à l'AR(1)}~: dans le cas d'un AR(1), la
loi marginale de l'observation initiale $y_1$ est simplement la loi
stationnaire univariée $\mathcal N(\mu,\gamma(0))$. Ici, il faut
écrire la loi \emph{jointe} de $(y_1,y_2)$, qui sous stationnarité est
une loi normale bivariée d'espérance $(\mu,\mu)'$ et de matrice de
variance-covariance~:
\[
\Sigma_2 = \begin{pmatrix} \gamma(0) & \gamma(1) \\ \gamma(1) & \gamma(0) \end{pmatrix},
\]
où $\gamma(0)$ et $\gamma(1)$ sont des fonctions \emph{non linéaires} des
paramètres $(\varphi_1,\varphi_2,\sigma_{\varepsilon}^2)$, déterminées
par les équations de Yule-Walker. C'est cette dépendance non linéaire
qui complique l'écriture (et l'optimisation numérique) de la
vraisemblance exacte.\newline

\question{2} La vraisemblance conditionnelle s'obtient en traitant
$(y_1,y_2)$ comme déterministes (« données ») plutôt qu'aléatoires~:
\[
\mathcal L_c(\bm\theta;\mathcal Y_T \mid y_1,y_2)
= \prod_{t=3}^T f(y_t\mid y_{t-1},y_{t-2};\bm\theta).
\]
Soit explicitement~:
\[
\mathcal L_c = \prod_{t=3}^T \frac{1}{\sqrt{2\pi\sigma_{\varepsilon}^2}}
e^{-\frac{(y_t - c - \varphi_1 y_{t-1} - \varphi_2 y_{t-2})^2}{2\sigma_{\varepsilon}^2}}.
\]
Elle diffère de la vraisemblance exacte par l'\emph{absence du
facteur} $f(y_1,y_2;\bm\theta)$. Son intérêt pratique est double~:
(i) elle ne fait intervenir les paramètres que de manière
\emph{linéaire} dans la moyenne conditionnelle, ce qui rend
l'optimisation beaucoup plus simple~; (ii) la perte d'information
($f(y_1,y_2)$) est négligeable lorsque $T$ est grand, et
l'estimateur du maximum de vraisemblance conditionnel reste convergent
et asymptotiquement équivalent à celui du maximum de vraisemblance
exacte.\newline

\question{3} La log-vraisemblance conditionnelle est~:
\[
\begin{split}
\ln \mathcal L_c =\;& -\frac{T-2}{2}\ln(2\pi\sigma_{\varepsilon}^2) \\
&- \frac{1}{2\sigma_{\varepsilon}^2}\sum_{t=3}^T (y_t - c - \varphi_1 y_{t-1} - \varphi_2 y_{t-2})^2.
\end{split}
\]
À $\sigma_{\varepsilon}^2$ fixé, seul le second terme dépend de
$(c,\varphi_1,\varphi_2)$~; sa maximisation revient à \emph{minimiser}
la somme des carrés des résidus~:
\[
S(c,\varphi_1,\varphi_2) = \sum_{t=3}^T (y_t - c - \varphi_1 y_{t-1} - \varphi_2 y_{t-2})^2.
\]
C'est exactement le critère des moindres carrés ordinaires appliqué à
la régression de $y_t$ sur la constante, $y_{t-1}$ et $y_{t-2}$, à
partir de $t=3$. En posant~:
\[
\bm{y} = \begin{pmatrix} y_3 \\ y_4 \\ \vdots \\ y_T \end{pmatrix},\quad
\bm{X} = \begin{pmatrix} 1 & y_2 & y_1 \\ 1 & y_3 & y_2 \\ \vdots & \vdots & \vdots \\ 1 & y_{T-1} & y_{T-2} \end{pmatrix},
\]
$\bm{X}$ étant de dimension $(T-2)\times 3$, l'estimateur s'écrit~:
\[
\widehat{\bm\beta} = \begin{pmatrix} \hat c \\ \hat \varphi_1 \\ \hat \varphi_2 \end{pmatrix}
= (\bm{X}'\bm{X})^{-1}\bm{X}'\bm{y}.
\]
On retrouve enfin l'estimateur de $\sigma_{\varepsilon}^2$ en
remplaçant les paramètres par leurs estimations dans la condition de
premier ordre par rapport à $\sigma_{\varepsilon}^2$~:
\[
\widehat{\sigma_{\varepsilon}^2} = \frac{1}{T-2}\sum_{t=3}^T \left(y_t - \hat c - \hat\varphi_1 y_{t-1} - \hat\varphi_2 y_{t-2}\right)^2.
\]

\bigskip

\exercice{4}\newline

\question{1} L'équation du processus latent peut s'écrire $(1-\rho L)
x_t = \eta_t$. En appliquant l'opérateur $(1-\rho L)$ à l'équation de
mesure $y_t = x_t + \varepsilon_t$, on obtient~:
\[
(1-\rho L) y_t = (1-\rho L) x_t + (1-\rho L)\varepsilon_t,
\]
soit~:
\[
(1-\rho L) y_t = \eta_t + \varepsilon_t - \rho\varepsilon_{t-1}.
\]
Notons $\xi_t \equiv \eta_t + \varepsilon_t - \rho\varepsilon_{t-1}$
le membre de droite.\newline

\question{2} Calculons les autocovariances de $\xi_t$. Comme
$\eta$ et $\varepsilon$ sont des bruits blancs indépendants entre eux
et à toutes les dates~:
\[
\gamma_{\xi}(0) = \mathrm{Var}(\eta_t) + \mathrm{Var}(\varepsilon_t) + \rho^2 \mathrm{Var}(\varepsilon_{t-1})
= \sigma_{\eta}^2 + (1+\rho^2)\sigma_{\varepsilon}^2.
\]
\[
\gamma_{\xi}(1) = \mathbb E[(\eta_t+\varepsilon_t-\rho\varepsilon_{t-1})(\eta_{t-1}+\varepsilon_{t-1}-\rho\varepsilon_{t-2})]
= -\rho\sigma_{\varepsilon}^2.
\]
Pour $h\geq 2$, tous les croisements font intervenir des innovations à
des dates distinctes, donc indépendantes~: $\gamma_{\xi}(h) = 0$.\newline

La fonction d'autocovariance de $\xi_t$ est donc identique à celle
d'un MA(1)~: $\gamma_{\xi}(h) = 0$ pour $h\geq 2$, et $\gamma_{\xi}(0)
\neq 0$, $\gamma_{\xi}(1)\neq 0$ (pour $\rho\neq 0$). On peut donc
écrire $\xi_t = u_t - \theta u_{t-1}$ avec $u_t$ un bruit blanc de
variance $\sigma_u^2$, où $\theta$ et $\sigma_u^2$ sont déterminés par~:
\[
\begin{cases}
(1+\theta^2)\sigma_u^2 = \sigma_{\eta}^2 + (1+\rho^2)\sigma_{\varepsilon}^2,\\
-\theta\,\sigma_u^2 = -\rho\sigma_{\varepsilon}^2,
\end{cases}
\]
soit~:
\[
\boxed{\;
\begin{cases}
(1+\theta^2)\sigma_u^2 = \sigma_{\eta}^2 + (1+\rho^2)\sigma_{\varepsilon}^2,\\
\theta\,\sigma_u^2 = \rho\,\sigma_{\varepsilon}^2.
\end{cases}\;}
\]
La condition d'inversibilité $|\theta|<1$ permet de sélectionner
l'unique solution acceptable.\newline

\question{3} En substituant la représentation MA(1) de $\xi_t$ dans
l'équation $(1-\rho L)y_t = \xi_t$, on obtient~:
\[
y_t = \rho y_{t-1} + u_t - \theta u_{t-1},
\]
ce qui est bien un processus ARMA(1,1) de paramètre auto-régressif
$\rho$, de paramètre moyenne mobile $\theta$ et d'innovation $u_t$ de
variance $\sigma_u^2$, ces deux dernières quantités étant déterminées
par le système ci-dessus.\newline

\question{4} La racine du polynôme retard auto-régressif est $1/\rho$,
strictement supérieure à 1 en module sous l'hypothèse $|\rho|<1$~: le
processus $y_t$ est asymptotiquement stationnaire au second ordre. Par
construction (voir la question 2), $\theta$ peut être choisi tel que
$|\theta|<1$, ce qui assure également l'inversibilité de la
représentation ARMA(1,1).

\bigskip

\textit{Remarque.} Cet exercice illustre une situation très courante~:
\emph{une dynamique AR(1) observée avec une erreur de mesure devient
  un ARMA(1,1)}. Plus généralement, l'observation bruitée d'un AR($p$)
donne un ARMA($p$,$p$). Ce résultat motive l'usage des
représentations état-mesure et du filtre de Kalman, qui permettent de
travailler directement avec le modèle d'origine plutôt qu'avec sa
forme ARMA réduite.

\end{document}

%%% Local Variables:
%%% mode: latex
%%% TeX-master: t
%%% End:
